\section{How the Calendar Works}

\subsection{Three calendars and one occurrence}

The Proper of Time supplies Sundays, seasons, and celebrations whose dates are fixed by Christmas or Easter. The General Roman Calendar supplies the universal fixed cycle. The United States proper adds or changes only the rows expressly approved for that territory. On a concrete date those layers are joined with the diocesan and local proper, then resolved by the table of liturgical days and the rules of occurrence.

\subsection{Ranks used in this reference}

\begin{longtable}{>{\bfseries\raggedright\arraybackslash}p{0.19\linewidth}>{\raggedright\arraybackslash}p{0.73\linewidth}}
\toprule
Rank or status & Meaning in the inventory\\
\midrule
\endhead
Solemnity & Highest ordinary rank; normally has First Vespers and proper presidential texts. Precedence still depends on the table of liturgical days.\\
Feast & Celebrated within the day; displaces an ordinary weekday but yields to higher days.\\
Memorial & Obligatory unless displaced or reduced under the norms. In Advent from 17--24 December, the Christmas octave, and Lent, a memorial occurring on a weekday may be observed only in the restricted manner allowed for a commemoration.\\
Optional memorial & May be celebrated when the day and pastoral judgment permit. If several optional memorials are listed, one may be chosen; they are not cumulative obligations.\\
Restricted commemoration & The limited observance of a memorial permitted by the norms on certain privileged weekdays; not a separate general-calendar rank parallel to a memorial.\\
All Souls commemoration & The sui generis observance on 2 November, placed at position 3 in the table of liturgical days. It is not the restricted memorial mechanism and precedes an Ordinary Time Sunday.\\
Required day of prayer & A nationally inscribed observance carrying a liturgical requirement without becoming a solemnity, feast, or memorial; the day actually occurring still governs the permitted formulary.\\
Optional special day & A nationally supplied civic or pastoral option whose use depends on the Missal's permissions and the day actually occurring; it does not displace a higher day.\\
\bottomrule
\end{longtable}

All unqualified entries in the 2002 Latin calendar are optional memorials by that calendar's own note. The tables therefore write the rank out rather than relying on typography. The rank belongs to the celebration, not to every possible local observance of the same saint.

\subsection{Precedence, transfer, and omission}

Sunday is the primordial feast day, but not every Sunday has the same precedence. Sundays of Advent, Lent, and Easter outrank all solemnities; a solemnity impeded by one of them is normally transferred under the norms. Sundays in Christmas Time and Ordinary Time yield to solemnities and feasts of the Lord listed in the General Calendar. Lower celebrations that meet a higher day are ordinarily omitted for that year, not moved at will.

The days of Holy Week from Monday through Thursday, the days within the Easter octave, and the weekdays of Advent from 17 through 24 December have special precedence. Particular-calendar solemnities and feasts must be added before an actual Ordo is resolved. The national transfer rules stated later are approvals, not examples of private discretion.
